\chapter{Layout as a Design Tool: MigLayout First}
\label{ch:miglayout}
\chapterquote{A layout manager is not merely a way to avoid coordinates; it is the written geometry contract of the screen.}{The layout maxim}

\begin{chaptermap}
Core question & How should ordinary production Swing screens be laid out today?\\
Primary APIs & \api{LayoutManager2}, \api{JPanel}, \api{JScrollPane}, \api{JLabel}, \api{JTextField}, \api{JButton}.\\
Primary library & \miglayout{} for dense, maintainable forms, dialogs, filter bars, and workspaces.\\
Prerequisite & Chapter~\ref{ch:components}: preferred size, validation, repainting, focus, scroll panes, and component geometry.\\
Companion example & \migLayoutExample.\\
\end{chaptermap}

\section{Why MigLayout is the default}

Swing's built-in layout managers are useful foundations. \api{BorderLayout} teaches regions. \api{FlowLayout} teaches preferred sizes. \api{BoxLayout} teaches axes. \api{GridLayout} teaches uniform cells. \api{GridBagLayout} teaches constraints, weight, fill, anchors, and insets. None of these lessons should be forgotten.

But most production screens are not pure examples of one lesson. A business dialog wants labels aligned by baseline, fields that grow horizontally, related and unrelated gaps, command buttons right-aligned at the bottom, optional sections, scrollable form bodies, and occasionally a table or message area that consumes the remaining height. \miglayout{} is a practical language for that geometry.

The library is not magic. It asks the same questions described in Chapter~\ref{ch:components}: what is each component's preferred size, minimum size, maximum size, baseline, and current visibility? Then it combines those answers with panel, column, row, and component constraints. It is powerful because it makes the parent-side contract visible.

The most useful way to read this chapter is therefore as the other half of the
component contract. Chapter~\ref{ch:components} explained how a child reports
natural size, baseline, opacity, focus, input behavior, and scrollable extent.
MigLayout explains how the parent uses that evidence. When a layout fails, do
not ask "what MigLayout incantation is missing?" first. Ask whether the child
reported truthful geometry and whether the parent constraints express the
desired relationship.

\begin{principlebox}{Layout text is code}
A MigLayout constraint string is part of the source contract. Keep it short, local, and reviewed. A good constraint says what grows, what aligns, where related fields live, and which region absorbs extra space.
\end{principlebox}

The most common mistake with MigLayout is to treat it as a shorthand for pixel
placement. That misses its real value. MigLayout is a compact notation for
relationships: this label belongs to this field, this field grows before that
button, this table owns extra height, this command row remains near the
decision point, this optional area collapses when hidden, and this form behaves
as a vertical document inside a viewport. A constraint string should read like
that relationship, not like the transcript of a visual tug-of-war.

There is a useful analogy with types. A Java type does not merely satisfy the
compiler; it documents what values a method accepts. A MigLayout constraint does
not merely satisfy the layout manager; it documents what geometry a panel
accepts. When the panel is changed later, the constraint should still tell the
maintainer which part of the screen is allowed to absorb space and which part is
not. If the maintainer must run the application after every small edit to
understand the intent, the layout is underdocumented even if it renders.

This is why examples in the book keep constraint strings near the panels they
describe. A shared constant such as \api{FORM\_COLUMNS} can be useful when a
team intentionally standardizes form geometry. A global constant named
\api{DEFAULT\_LAYOUT} is less helpful if it hides the interesting decision from
the reader. Prefer local constraints until a repeated shape has proven itself.
Then extract the repeated shape with a name that explains the geometry, not with
a name that merely says "standard".

MigLayout also rewards reading the screen from large regions to small
relationships. First identify the shell: header, body, status, command row,
sidebar, table, detail, preview. Then identify each region's local grammar:
label/editor rows, dense filters, button clusters, tables, notes areas, warning
strips. Finally write component constraints. If the component constraints are
written before the regions are understood, they often become a long list of
special cases.

\textbf{The three constraint layers.}
MigLayout has three constraint layers. The constructor usually receives layout constraints, column constraints, and row constraints. Each component added to the panel may also receive component constraints.

\begin{lstlisting}[caption={The three MigLayout constructor constraints plus component constraints.}]
JPanel form = new JPanel(new MigLayout(
        "fillx, insets 16, gap 8",  // layout constraints
        "[][grow,fill]",            // column constraints
        "[][][]push[]"));           // row constraints

form.add(new JLabel("Name"));
form.add(nameField, "wrap");        // component constraints
form.add(new JLabel("Type"));
form.add(typeBox, "wrap");
form.add(buttons, "span 2, align right");
\end{lstlisting}

The layout constraints describe the panel as a whole. The column constraints describe vertical tracks. The row constraints describe horizontal tracks. The component constraints describe how one child participates in the grid.

This is why long, clever component constraints are often a smell. If every component says \api{growx}, \api{pushx}, \api{span}, \api{split}, \api{gapleft}, and \api{align}, the parent geometry may not be well described. Prefer a clear column that grows and fills over repeating the same behavior on every field.

The layers should not fight each other. If a column says that editors grow and
fill, each editor should not need to repeat the same instruction. If a row is
the only vertical consumer, the components in that row should not also carry
unrelated push instructions merely to make the screenshot look right. The more a
child constraint repeats a parent rule, the harder it is to discover which rule
actually matters. Repetition in layout strings is not harmless; it is competing
documentation.

The component constraint is best used for facts that are local to that
component's participation in the grid. A field may wrap because it ends a row. A
button panel may span both columns because it belongs to the whole form. A help
link may split a cell with another small control. A warning label may span the
editor column but not the label column. These facts cannot be stated cleanly in
column or row constraints because they are about one child.

Column and row constraints should describe tracks, not widgets. A column can be
the label column, the editor column, the action column, or the gutter. A row can
be a preferred field row, a growing notes row, a warning row, or a command row.
Once a row constraint is meaningful only because "the third component happens to
be there", the layout is too dependent on add order. MigLayout still uses add
order, of course, but the design should be visible at the grid level.

The constructor form is not the only readable form. For large panels, teams
often build small helper methods for region construction: \api{createFilterBar},
\api{createEditorBody}, \api{createCommandRow}. The important part is that each
helper owns one coherent geometry. A single method that creates every component
and adds them all to one huge grid can be as hard to read as \api{GridBagLayout}
code. MigLayout reduces ceremony; it does not remove the need for composition.

\section{Vocabulary before notation}

MigLayout is easiest to learn when its words are understood before its
abbreviations are memorized. A constraint such as \api{span 2, growx, wrap}
looks cryptic only if the reader has not yet formed the picture behind it. The
picture is simple: the parent has a grid; rows and columns are tracks in that
grid; a component starts in a cell; the component may occupy more than one
cell; extra space is assigned according to growth rules; the component may or
may not fill the space it receives; gaps express visual relationship; wrapping
moves the insertion point to the next row.

This vocabulary is not ornamental. It is the language used when debugging a
screen. "The field does not grow" is an ambiguous complaint. "The editor column
does not grow" is one complaint. "The editor column grows, but the field does
not fill" is another. "The field fills, but its maximum size prevents useful
width" is a third. A team that knows the vocabulary can discuss layout without
inventing a new story for every screenshot.

Figure~\ref{fig:miglayout-vocabulary} shows the mental model. It is not a
complete MigLayout manual, but it names the relationships that appear in nearly
every useful panel. The vocabulary should be read together with the previous
chapter's component contract: MigLayout arranges components by asking them for
minimum, preferred, maximum, and baseline information, then applying parent
constraints to the available space.

\begin{figure}[htbp]
\centering
\includegraphics[width=0.96\linewidth]{\migLayoutVocabularyFigure}
\caption{MigLayout vocabulary as geometry. The component begins in a cell; columns and rows form tracks; a component may span cells; gaps separate related and unrelated facts; grow receives slack; fill consumes the allocated cell; wrap moves the next component to the following row.}
\label{fig:miglayout-vocabulary}
\end{figure}

Read Figure~\ref{fig:miglayout-vocabulary} as a sequence rather than as a
poster. First find the grid. The grid belongs to the parent panel, not to any
one child. Then find the starting cell for a component. The starting cell is the
place where the component enters the grid. Next look for span: a component can
occupy more than its starting cell when the visual relationship is wider than
one track. Only after those facts are clear should the reader think about grow,
fill, and wrap. This order mirrors the debugging order: where is the component,
what space is available, what space was allocated, and how did the component
use it?

The figure also separates relationship from measurement. Gaps are shown as
meaningful spaces between related and unrelated controls; they are not merely
empty pixels. Grow is shown as slack distribution; it is not the same as a text
field's preferred width. Fill is shown as a component using the cell it has
received; it is not the same as receiving extra space. Wrap is shown as the
insertion point moving to the next row; it is not a visual line break inside a
component. These distinctions are small enough to learn quickly and important
enough to prevent many wrong fixes.

The vocabulary is intentionally larger than a short glossary because layout
mistakes usually arise between terms. A component may be in the right cell but
not span enough columns. It may span correctly but sit in a non-growing column.
The column may grow while the component does not fill. The component may fill
but report a maximum size that prevents useful expansion. A form may grow
correctly until it is placed inside a scroll pane whose view does not track
viewport width. Once the terms are known separately, these combined failures can
be named instead of guessed.

\textbf{A cell is a starting position, not a widget.} A cell is the grid
location at which a component begins. The component is then allowed to occupy
that one cell or to span additional cells. This distinction matters because the
grid is a parent-side structure, while the component is a child-side object. A
text field is not "the second column"; it is a component whose starting cell is
in the editor column for one row. Another field may start in the same column on
the next row. A validation summary may start in the editor column and span the
command column. The cells describe the parent grammar; the components supply
content.

The common beginner mistake is to think that each visual object creates a new
permanent column. A form with a unit label, a help button, and a required
marker can quickly become a seven-column grid if every local detail is promoted
to a global track. That is usually the wrong abstraction. The main grid should
describe repeated structure: labels, editors, action column, perhaps a marker
column. Local composition belongs in a split cell or a nested panel when the
exception is truly local. The goal is not to minimize the number of columns.
The goal is to make the grid's columns meaningful across the region.

A cell also explains why add order remains important. MigLayout is not a
declarative constraint solver that ignores insertion order. It places
components as they are added, subject to the current cell, wrapping, spanning,
splitting, and other constraints. Good source code therefore makes add order
follow reading order. If a user reads label, field, help, warning, command, the
source should not add those objects in a surprising order merely because the
constraints can still produce the desired picture. Layout source is a reading
artifact as well as a rendering artifact.

\textbf{A track is the row or column policy.} The word "track" is useful
because rows and columns are not merely locations. They are policies for a
strip of space. A label column may be preferred width. An editor column may
grow and fill. A command column may stay at preferred width. A notes row may be
the only vertical consumer. A warning row may be preferred height but may
appear and disappear. When the chapter says "column constraint" or "row
constraint", it is talking about those track policies.

Track thinking prevents a familiar maintenance failure. Suppose every text
field in a form is added with \api{growx, pushx}. The screen may look correct,
but the source hides the actual rule: the editor column is the thing that
should receive extra width. If a later developer adds a combo box and forgets
the local growth flags, that one row behaves differently. If the editor column
itself says \api{grow, fill}, the repeated policy has one home. The component
constraint can then say only what is exceptional about that component's
placement.

Rows deserve equal discipline. A panel that grows vertically should say which
row receives the extra height. A table row, notes area, preview pane, or empty
push row can legitimately be the vertical consumer. They are different design
answers. If every row grows, the form may stretch fields into awkward heights.
If no row grows, a tall window may produce unclaimed whitespace or strange
alignment. A row constraint is therefore not housekeeping; it states which
part of the screen benefits from vertical space.

\textbf{Grow allocates extra space.} Growth begins when the parent has more
space than the preferred layout requires. MigLayout must decide where that
slack goes. A growing column may become wider. A growing row may become taller.
A growing component may receive more space inside the cell or cells it
occupies. In the common form case, the editor column grows horizontally, while
the label column remains at preferred width. The result is a form that becomes
more useful as the window widens: the fields gain editing room instead of the
labels drifting away from them.

Growth is a negotiation, not a command that one object executes in isolation.
If the panel itself is not given extra width by its parent, the editor column
cannot grow. If the column grows but the component does not fill, the field may
still sit at preferred width inside a larger cell. If the component has a
restrictive maximum size, it may refuse to become wider. This is why layout
debugging must inspect containment, track constraints, component constraints,
and component size contracts together. A single word in one layer is rarely the
whole story.

Growth should also have priority. A filter bar may contain a search field, a
status combo box, a date field, and a Search button. It is usually wrong to let
all of them grow equally. The search field may deserve the slack because it
contains arbitrary text. The status combo box has a bounded vocabulary. The
date field has a predictable length. The button is a command, not a data
region. Writing a growth policy is therefore an act of interface design: it
tells the application what kind of user work benefits from extra space.

\textbf{Fill uses the space that has already been assigned.} Fill is often
confused with grow because good form fields frequently need both. They are
nevertheless different ideas. Grow is about receiving extra space from the
layout. Fill is about occupying the space once it has been received. A field in
a wide cell may remain at preferred width if it does not fill. A button in a
wide cell might deliberately avoid filling because stretched command buttons
can look clumsy. A scroll pane, text field, table, and notes area commonly
fill. An icon, button, color swatch, or compact status label commonly does not.

This distinction gives a useful diagnostic sentence: "The cell is large, but
the component is not using it." That sentence points toward fill, alignment,
or maximum size. It does not point toward the column's growth policy. Without
the distinction, teams often add more \api{growx} declarations and become
frustrated when nothing changes. The field already received space; it simply
did not fill it.

Fill also interacts with baseline and alignment. A component can fill
horizontally and still align its text by baseline with neighboring labels. A
component can fill vertically and make a row look wrong if it was meant to stay
at preferred height. This is why form rows usually grow horizontally but not
vertically. Width helps text entry. Arbitrary height often just makes controls
look bloated.

\textbf{Shrink is the other half of grow.} Many layout explanations discuss
only what happens when the window becomes larger. Production screens also
become smaller: a laptop display, a remote session, a split workspace, a user
font increase, a translated resource, or a docked tool window can all compress
the panel. The layout must decide what gives way. Does the search field become
shorter? Does the table scroll horizontally? Does a label wrap? Does an
advanced filter move to another line? These are not merely visual details; they
determine whether the screen remains usable under ordinary conditions.

MigLayout constraints can express minimum and maximum sizes, but the first
question is design. A fixed-width currency code may shrink little because its
domain is short. A customer name field may shrink but should keep a sensible
minimum. A table may keep its height claim and let columns manage their own
compression. A command row should usually remain visible and preferred height.
When a panel has no shrink story, it often survives the developer's monitor and
fails everywhere else.

The shrink story should be tested deliberately. Open the example at a smaller
size. Increase the default font. Replace labels with longer resource strings.
Use FlatLaf light and dark themes. Put the form inside a scroll pane. If the
screen fails, do not immediately add magic widths. Ask which track is allowed
to surrender space and which component's minimum size represents a real
usability promise.

\textbf{Span means one component occupies several cells.} A component has one
starting cell, but it may span horizontally, vertically, or both. The common
form examples are command rows that span the label and editor columns, warning
messages that span the editor area, and scroll panes that span several columns
in a workspace. Spanning is a statement that the component belongs to a larger
region than the ordinary field row.

Span becomes messy when it is used to compensate for a poor grid. If half the
components span two columns and the other half span three, the grid may not be
describing the screen's repeated structure. Sometimes that is legitimate: a
workspace is less regular than a form. Often it means the region should be
split into smaller panels with different local grammars. A filter bar, result
table, detail form, and status strip do not need one universal grid.

Spanning should preserve reading order. A validation summary that spans the
editor column and command column should still appear near the fields it
explains. A section heading that spans all columns should start a group, not
float between unrelated controls. A button row that spans the form should
communicate that the commands apply to the whole form. When span is merely a
way to push pixels into place, the source becomes harder to review than the
geometry it tries to express.

\textbf{Split is local composition inside a cell.} A split cell allows more
than one component to participate in the same cell. This is useful for small
local relationships: a text field with a unit label, a search field with a
clear button, a compact pair of date fields, or a group of tiny command
buttons. Split should not be used to build an entire toolbar inside a form
cell unless that toolbar is genuinely part of the field's local interaction.

The difference between span and split is worth keeping clear. Span says one
component occupies more of the parent grid. Split says one parent-grid position
contains a small local sequence. If the local sequence has its own growth,
alignment, visibility, and lifecycle rules, a nested panel may be clearer than
a long split constraint. MigLayout makes split convenient; it does not remove
the need to name boundaries.

Split is also where accessibility and focus order can become subtle. If a
field cell contains a text field, clear button, and help link, keyboard
navigation should still make sense. The source order should match the intended
focus order. Accessible names should explain the small controls. Corusco can
provide stable help topics and commands, but the view still decides that these
controls are local companions of one field rather than global actions.

\textbf{Wrap controls the flow of additions.} A wrap says that the next
component begins on the next row. In a two-column form, a row is often written
as label, field with \api{wrap}. That is readable because it mirrors the
screen. In denser panels, explicit wrapping prevents a later addition from
silently continuing the previous row. It also documents which component ends
the current thought.

Wrap is a small word with large review value. A missing wrap can make a new
field appear in the wrong row. An accidental extra wrap can create a blank
line or a jagged form. When reviewing a form, read each row as a sentence. If
the sentence ends after the editor, wrap there. If a local help button belongs
to the same row, do not wrap too early. If a section heading starts a new
group, wrap before it and consider an unrelated gap.

It is tempting to treat wrap as punctuation that can be fixed by running the
program. Prefer to make it visible in the source. A line ending with
\api{"wrap"} is often easier to review than a clever default wrap count hidden
in the layout constraints. Defaults are useful for uniform grids, but ordinary
forms usually benefit from explicit row endings, especially while the screen is
still evolving.

\textbf{Gap is visual grammar.} A gap is the space between components or around
panel content. The important point is not the numeric size. The important point
is the relationship it expresses. Related controls are closer. Unrelated groups
are farther apart. Insets keep content from touching the panel boundary.
Command rows often receive separation from the editable body. A warning may sit
near the field it explains, while a form-level summary may sit above the group
it governs.

Random gaps make a screen look unsettled. They also make maintenance harder
because a later developer cannot tell whether a particular distance is
semantic, accidental, or a compensation for another layout defect. Use
platform gaps, named gap policies, or a small local vocabulary of spacing
decisions. If a gap is unusual, a short comment may be justified. If every gap
requires a comment, the region may be visually overcomplicated.

Gaps are especially important when screenshots are used as documentation. A
single screenshot can hide whether spacing survives font changes. A layout
whose gaps express relationships will usually remain readable as metrics
change. A layout whose gaps are tuned as pixels will not. This is why the book
examples prefer modest, repeated gaps and let component preferred sizes do
their work.

\textbf{Insets are the panel's boundary grammar.} Insets are gaps between the
panel edge and its content. They are easy to overuse because nested panels can
each add their own insets. A dialog body may need insets. A nested panel that
already sits inside a padded region may not. A card-like repeated item may need
internal padding because it is a visual object. A transparent grouping panel
inside a form may need no extra insets at all.

When a screen feels both cramped and wasteful, inspect nested insets. The
content may be squeezed by several layers of padding even though the outer
window has plenty of space. This is not a MigLayout-specific failure; it is a
general Swing containment failure. MigLayout simply makes the policy visible
enough to review. The rule is to decide which boundary owns the breathing room.
Do not let every boundary claim it by habit.

\textbf{Push claims unassigned slack.} Push is often used to move a region
away from another region or to make one area claim leftover space. A common
form pattern uses a push row before the command row, so the buttons remain near
the bottom while the fields keep preferred height. A status strip may use push
to separate left-side messages from right-side counters. A workspace may use
push so the central table receives the height not needed by filters and status.

Push should be used with intent. It is not a general remedy for "the component
is in the wrong place." Ask what slack exists and who should own it. If a form
body is a vertical document, the slack may belong to the scroll pane rather
than to a blank row. If a detail pane has a notes area, the notes row may be a
better owner than an empty push. If the command row should remain adjacent to
the final field, a bottom push may be wrong. Push is a spatial policy, not a
decoration.

Push also makes invisible policy visible. Without it, extra space may be
distributed in a way that looks acceptable on one size but strange on another.
With it, the source says where unused space belongs. This is one of the reasons
MigLayout is valuable in business screens: ordinary layout code can state
ordinary workflow priorities without a page of \api{GridBagConstraints}.

\textbf{Alignment describes where a component sits inside its allocation.}
After size has been assigned, a component still needs a position inside the
cell or span. It may align leading, trailing, centered, baseline, top, or
bottom depending on the context. Alignment is not the same as fill. A component
that fills has little remaining internal positioning choice. A component that
does not fill must still be placed.

Alignment is where many screens reveal whether they were designed or merely
assembled. Text fields usually align as a column. Labels usually align with
their fields by baseline and often by trailing edge in left-to-right forms.
Small buttons may align with the field they supplement. Command rows usually
align according to platform and workflow convention. Icons and color swatches
may align center. A screen with arbitrary center and right alignment reads
poorly even if all components are technically visible.

Use leading and trailing language mentally even when a constraint string uses a
concrete word. Localization and right-to-left orientation make absolute
"left" and "right" assumptions less portable. The application may still choose
a platform-specific command order, but the design idea should be explicit:
commands are trailing, Help may be separated, destructive commands may be
visually distinct, and field labels follow the reading direction.

\textbf{Baseline is typographic alignment.} Swing components can report a
baseline, the line on which their text sits. MigLayout can use that information
to align labels, fields, combo boxes, and buttons so that a row reads as one
line of text rather than as a stack of boxes. Baseline alignment is one of the
small differences between a professional form and an amateur arrangement.

Baseline depends on the child component. Standard labels, fields, buttons, and
combo boxes usually provide useful baselines. A custom component may not. A
custom date editor, token field, formatted field, or search component that is
intended to sit in a text row should decide what its baseline is. Otherwise the
parent layout must fall back to top, center, or bottom alignment, and the row
may look slightly wrong at every font size.

Validation decoration can disturb baseline if it changes borders or insets
while the user types. A red border that adds pixels around the field may change
preferred size and baseline. A problem icon inserted into the row may push the
editor. A stable decoration area, overlay, or carefully composed border keeps
the row from jumping. The layout vocabulary therefore meets the validation
vocabulary: a problem is semantic, but its visible mark still participates in
geometry.

\textbf{Size groups express equality when equality matters.} Some screens need
several controls to share a width even though their preferred sizes differ.
Button rows are a common example: OK, Cancel, and Apply may use equal widths
under a chosen platform policy. A pair of short numeric fields may have the
same width because the domain says they are comparable. Size grouping should
not be a reflex. Equal widths can clarify peer controls; they can also waste
space or imply a relationship that does not exist.

The useful test is semantic. If two controls are peers in the user's decision,
a shared size may help. If one field accepts a name and another accepts a
postal code, equal width is probably less important than giving the name field
room to breathe. MigLayout can state equality, but the screen must justify it.
This mirrors Corusco's broader design: descriptors and keys can make identities
stable, but visual grouping remains a presentation decision.

\textbf{Docking, absolute positions, and bounds are escape hatches.} MigLayout
has features for special cases, but ordinary forms should not reach for
absolute coordinates. The moment a panel depends on specific pixel positions,
it becomes fragile under font, Look-and-Feel, localization, and scaling
changes. A rich visual component such as a map or canvas may use coordinates
inside its own painting model. The parent panel should still place that
component through normal layout relationships.

Escape hatches are not forbidden. They are expensive. Use them when the screen
really has a spatial model that cannot be described as rows, columns, tracks,
and regions. Then isolate the special component. Do not let one unusual visual
region infect the surrounding form with fixed coordinates. A GIS viewport,
paint editor canvas, or diagram component may be custom and coordinate-heavy;
the toolbar, properties pane, status strip, and dialogs around it can still be
ordinary MigLayout panels.

\textbf{Debug mode shows geometry, not intent.} MigLayout can draw layout
debugging marks. Temporary borders and background colors can also reveal
regions. These tools are useful because they make cells, gaps, and spans
visible. They do not answer the design question. A debug grid can show that a
notes row receives extra height. It cannot tell whether the notes row should
receive that height. A debug grid can show that a warning spans three columns.
It cannot tell whether the warning belongs near the field or at the form
summary.

Debugging therefore uses two passes. First identify the mechanical fact: which
cell, which track, which span, which growth rule, which fill rule, which gap,
which component size. Then identify the design fact: which region is primary,
which relationship is being expressed, which information should remain visible,
which model fact drives a visual state. The mechanical pass explains what
happened. The design pass decides what should happen.

\textbf{Corusco names meaning; MigLayout names geometry.} A generated Corusco
form can name a field, expose a typed model, bind raw text to semantic state,
attach validation problems, define command identity, and provide help topics.
None of that tells the view where the field should appear. MigLayout describes
the geometry: label column, editor column, growing notes row, trailing command
row, fixed status strip, scrollable body. This division is a strength. It lets
the same stable field identity appear in a compact dialog, a wide workspace
detail pane, or a wizard page without pretending that all three screens share
one visual grammar.

The division also improves review. If a Save button is enabled incorrectly,
inspect command state, form committability, bindings, and validation problems.
If the Save button is in the wrong place, inspect MigLayout constraints and
platform command policy. If a field shows the wrong problem, inspect Corusco
problem targets and binding behavior. If the problem mark shifts the row,
inspect layout and decoration. A clean vocabulary keeps failures from becoming
one undifferentiated "the UI is wrong" complaint.

\textbf{Read a constraint string as a short design proof.} A useful MigLayout
review does not begin by asking whether the screenshot looks acceptable today.
It begins by asking whether the constraints prove the screen's intended
behavior. A two-column form should be explainable in a few sentences: the first
column is the label track, the second column is the editor track, the editor
track grows and fills, ordinary rows stay at preferred height, the notes row or
scroll pane owns vertical slack, and the command row spans the form. If those
sentences cannot be spoken, the constraint string is probably doing work that
is not yet understood.

This proof style is especially valuable for code review because it gives the
reviewer something better than taste. "The field looks too short" is a weak
review comment. "The editor column has no growth policy, so the field cannot
benefit from a wider dialog" is a precise comment. "The warning spans all
columns, so it reads like a form-level problem even though its target is one
field" is a precise comment. "The command row sits inside the scrollable body,
so OK may disappear while correcting long forms" is a precise comment. The
vocabulary turns visual unease into maintainable engineering observations.

The same habit helps when translating from \api{GridBagLayout}. Many old Swing
screens already contain the right ideas but express them with too much
ceremony: weight, fill, anchor, insets, grid width, and grid height are scattered
across mutable constraint objects. When migrating to MigLayout, do not
mechanically translate every line. First recover the design proof. Which column
was meant to grow? Which row held the main work surface? Which components were
peers? Which gaps expressed grouping? Then write the MigLayout version from the
proof, not from the accident of the old source.

\begin{detailbox}{Grow and fill are not synonyms}
A column with \api{grow} may become wider. A component with \api{growx} may be allocated extra width. A component with \api{fillx} uses the width it receives. In simple forms these often appear together, but confusing them makes debugging harder.
\end{detailbox}

Shrink deserves the same attention as grow. A screen is often designed on a
wide monitor, then used on a laptop, remote desktop, or split workspace. Which
component gives up space first? Which field may become narrow but still useful?
Which table should keep a minimum number of visible columns? Which label may
wrap? Which command row should remain visible? A layout that only states growth
policy is half specified; the shrinking behavior is where many production
annoyances appear.

The preferred-size chapter taught that a component's minimum size is partly a
usability promise. MigLayout lets the parent respect or override that promise,
but the team must choose. A search field may shrink to a modest width. A
currency field may have a strict minimum because partial numbers are dangerous.
A table may be allowed to lose horizontal space because it has a horizontal
scroll bar. A form with a required text field may prefer to keep the field
usable and let the window enforce a minimum size. These choices are design, not
syntax.

Alignment is also not a decoration added at the end. Leading and trailing
alignment affect how the screen reads. Baseline alignment affects whether the
form looks stable. Center alignment may be appropriate for icons or compact
tools but amateurish for text fields in a business form. Fill may be correct for
editors and tables but wrong for buttons. A good layout states these choices in
the smallest place where they are true.

\textbf{A form skeleton.}
Most form panels start with two columns: labels and editors. Labels normally use their preferred width. Editors normally fill a growing column. Rows normally use preferred height until one row is explicitly marked as the vertical consumer.

This pattern depends on the standard component size stories from
Chapter~\ref{ch:components}. Labels are allowed to be as wide as their text,
font, icon, and border require. Text fields are allowed to compute a natural
height from the Look-and-Feel and a natural width from their columns or content
policy. The layout then decides that the editor column, not the label column,
receives extra width. The result is robust because neither side lies.

\begin{lstlisting}[caption={A compact MigLayout form pattern.}]
JPanel form = new JPanel(new MigLayout(
        "fillx, insets 16, gap 8",
        "[][grow,fill]",
        "[][][]push[]"));
form.add(new JLabel("Name"));
form.add(nameField, "wrap");
form.add(new JLabel("Type"));
form.add(typeBox, "wrap");
form.add(new JLabel("Status"));
form.add(statusField, "wrap");
form.add(buttons, "span 2, align right");
\end{lstlisting}

This pattern is deliberately plain. The first column is only as wide as the labels require. The second column grows and fills. The \api{push} before the final row moves the button row toward the bottom if the parent becomes taller. The button row spans both columns and aligns right.

Plain does not mean simplistic. This two-column form pattern is the foundation
for a large amount of business UI because it follows the user's scanning path.
The label column names facts. The editor column carries values. The command row
is visually separated from data entry. The growing policy says that wider
windows should help editing, not merely add whitespace. The bottom push says
that extra height belongs to the body, not to random gaps between fields.

When the form grows beyond a few fields, group the rows by meaning before
adding visual chrome. Identity fields may form one group, contact fields
another, permissions another, notes another. Related fields may stay close.
Unrelated groups may receive a larger gap or a small heading. Avoid turning the
whole form into a stack of bordered boxes; most business forms need rhythm and
grouping, not card nesting. MigLayout can provide grouping through gaps,
spanning labels, and row constraints without making every group a visual panel.

A field that needs a unit label is still part of the editor relationship. For
example, a credit-limit field and a currency label may share the editor cell
with \api{split} or a small nested panel. The domain fact is one field with an
adjacent unit, not two unrelated columns in the entire form. Use the parent grid
for repeated structure and use a nested panel or split cell for local
composition. This keeps the main form columns from multiplying for one special
case.

A required marker should be handled with similar care. Some applications use a
small marker column between label and editor. Others include the marker in the
label text or use validation decoration. The choice should be consistent within
a screen. A random marker inserted as another component in an ordinary editor
cell will disturb alignment and reading order. If markers are layout elements,
give them a track. If they are decoration, let a behavior own them.

\begin{figure}[htbp]
\centering
\includegraphics[width=0.82\linewidth]{\migLayoutFigure}
\caption{A small MigLayout form from the companion examples. The field column grows, labels keep preferred width, and the command row is aligned right at the bottom.}
\label{fig:miglayout-form}
\end{figure}

\begin{figure}[htbp]
\centering
\includegraphics[width=0.96\linewidth]{\migLayoutDenseFigure}
\caption{MigLayout used for a denser business screen. The fields remain readable because growth, wrapping, related gaps, and the result table's vertical claim are stated in the layout constraints.}
\label{fig:miglayout-dense-screen}
\end{figure}

\textbf{Column and row thinking.}
A good MigLayout panel is often easier to read by columns and rows than by component additions. Ask first: which vertical tracks exist? A dense editor may have a label column, a required-marker column, a growing field column, and a small command column. A search bar may have a label, a growing search field, a fixed filter combo, and a command button.

Then ask which horizontal tracks exist. Most rows are preferred height. A table row or notes area may grow. A status row may sit at the bottom. A button row may remain preferred height but be separated from the body by unrelated gap.

When the grid is clear, component additions become simple. If the grid is unclear, component constraints become a pile of local corrections. Layout code should be read as design, not as a log of where individual widgets were shoved until the screenshot looked acceptable.

\textbf{Preferred sizes still matter.}
MigLayout begins with component size contracts. If a custom component reports a preferred width of 30 pixels, MigLayout will not magically infer that it is a chart needing 500 pixels. If a text field's maximum size prevents growth, a growing column may still not produce a growing field. If a scroll pane's preferred viewport size is too small, the layout may look cramped until constraints override it.

This is the first bridge from Chapter~\ref{ch:components}. Parent constraints cannot repair a lying child. They can only decide how to use the child's reported sizes. When layout behaves strangely, inspect both sides: the component's preferred/minimum/maximum sizes and the parent constraints.

\begin{pitfallbox}{fixing child lies in the parent}
It is possible to override sizes in MigLayout constraints, but doing so everywhere hides the real problem. If a reusable custom component always needs a better preferred size, fix the component. If one screen has a special visual policy, express that policy in the layout.
\end{pitfallbox}

\textbf{Gaps as grammar.}
Gaps are not blank pixels; they are grammar. A small related gap says two things belong together. A larger unrelated gap says a new group begins. Insets say the panel content should breathe inside its border. Button gaps should match platform and Look-and-Feel expectations where possible.

Hard-coded large gaps age poorly. They may look good at one font size and ridiculous at another. MigLayout can use logical gap names and platform-aware spacing. Even when using numeric gaps, keep them modest and consistent. A screen with random 3, 7, 11, 14, and 23 pixel gaps rarely looks intentional.

The grammar metaphor is useful during review. If a gap is increased, ask what
semantic separation it represents. If a gap is removed, ask whether two controls
now read as one unit. If a section receives a top gap, ask whether the preceding
section really ended. These questions keep spacing from becoming a matter of
personal taste. Taste still matters, but it is disciplined by the screen's
meaning.

Insets are similar. A panel inside a dialog usually needs insets because the
content should not touch the window border. A panel inside a cell of a larger
layout may need no insets because the parent already supplies rhythm. Nested
insets accumulate quickly. A screen can look padded and cramped at the same time
when every nested panel adds its own border insets. Inspect the nesting before
adding more space.

Use visual debugging sparingly but decisively. Temporary borders around regions
often reveal that the apparent gap is really a nested panel's inset, a scroll
pane border, a label's preferred width, or a row that grows unexpectedly. Remove
the temporary borders after the diagnosis. The final source should express
geometry through constraints, not through debugging artifacts.

\section{Baseline alignment}

Text in labels, fields, combo boxes, and buttons should usually align by baseline. Baseline alignment is one reason MigLayout forms look less amateur than hand-arranged panels. Components with different heights can still appear typographically aligned.

Baseline alignment depends on the components reporting meaningful baselines. Standard Swing components usually do. Custom components may not. If a custom editor appears slightly too high or low in a form row, inspect baseline behavior before adding arbitrary top gaps.

Baseline alignment is most visible when components have different heights. A
combo box, text field, button, and label may all occupy the same row but draw
their text at slightly different vertical positions. Top alignment makes the row
look accidental. Center alignment can look close but still wrong. Baseline
alignment makes the text read as one line. This small detail is one reason a
Swing form can feel professional without decorative styling.

Custom components that participate in forms should decide whether they are
text-like. A color swatch beside a label may not need a baseline; it can align
center. A custom date editor, token field, or search box probably should expose
a meaningful baseline. If implementing baseline behavior is too much for the
component, wrap it in a small panel whose alignment policy is explicit. Do not
hide the problem with arbitrary top gaps; those gaps will fail when the font
changes.

Baseline alignment also interacts with validation decoration. A red border or
problem icon should not push the field text off the baseline. Decoration that
changes border insets can change preferred size and baseline. A better policy
is to reserve decoration space, compose borders carefully, or paint markers in
an overlay. Layout and validation are separate concerns, but their visual
contracts meet in the row.

\textbf{Growing in one direction.}
Many screens should grow horizontally but not vertically. A search form might stretch fields when the window widens while keeping row heights stable. A notes field may grow in both directions. A table should usually grow in both directions. A command row should usually grow horizontally only as a container but keep buttons at preferred size.

The important habit is to name the growth consumer. If a panel has extra vertical space and no row claims it, the result may be awkward centering or bottom whitespace. If several rows claim it equally, the result may be stretched fields that should not stretch. One growing row is often clearer than many.

For forms, the vertical growth consumer is often a notes area, a table, a
preview, or an empty push row before the command row. Choose deliberately. If a
notes area grows, the user gets more editing space. If a table grows, the user
sees more rows. If an empty push row grows, the command row stays at the bottom
while the body remains at preferred height. Each choice communicates a workflow
priority.

Horizontal growth also needs priority. In a filter bar, the search field may
deserve most extra width while a status combo box remains preferred size. In a
details header, a title label may grow but an action button should not. In a
two-column editor, both editor columns may grow equally or one may dominate.
Avoid giving every component \api{growx} merely because the window can widen.
Growth without priority is just stretching.

When a screen is resized smaller, growth policy becomes compression policy.
MigLayout can only distribute the available space; it cannot decide which
business fact is more important unless the constraints say so. If the status
field must remain readable, give it a minimum. If the optional filter can
collapse or move into an advanced panel, make that a screen behavior. If a table
can scroll horizontally, accept that policy and keep the surrounding form
stable.

\textbf{Buttons and command rows.}
Button rows are not ordinary form rows. They represent decisions. In left-to-right desktop applications, primary and secondary commands often align to the right at the bottom of a dialog or form. Destructive commands may be separated. Help may sit on the opposite side. Platform conventions vary, but random placement is never a convention.

MigLayout can express command rows with \api{span}, \api{split}, \api{align right}, and push constraints. Corusco can supply the commands behind the buttons, but it should not decide the spatial convention. Layout remains part of application design.

\begin{coruscobox}{commands and layout}
Use Corusco commands to remove duplicated action state. Use MigLayout to place the buttons. A generated command descriptor should not become a layout policy.
\end{coruscobox}

\textbf{Scrollable form bodies.}
Long forms should usually scroll vertically and avoid horizontal scrolling. This is a contract among three objects: the scroll pane, the viewport view, and the form layout. The form panel should fill the viewport width, fields should grow within that width, and rows should determine the preferred height.

Swing does not provide a universal "width-tracking panel" out of the box. Some applications implement a small \api{Scrollable} panel whose \api{getScrollableTracksViewportWidth} returns true. Others use layout constraints and viewport listeners. The important point is conceptual: the view should track viewport width when the form is a vertical document of fields.

The scroll-pane discussion in Chapter~\ref{ch:components} gives the missing
child-side vocabulary: viewport extent is the visible window, the view preferred
size is the scrollable world, and \api{getScrollableTracksViewportWidth}
chooses whether the view accepts the viewport width. MigLayout can make the
fields grow inside the form, but it cannot by itself decide whether the form is
a width-tracking scrollable view. That decision belongs to the component placed
inside the viewport.

\begin{lstlisting}[caption={A small width-tracking form panel for scroll panes.}]
public final class ScrollableFormPanel extends JPanel implements Scrollable {
    public ScrollableFormPanel() {
        super(new MigLayout("fillx, insets 16", "[][grow,fill]", ""));
    }

    @Override
    public boolean getScrollableTracksViewportWidth() {
        return true;
    }

    @Override
    public boolean getScrollableTracksViewportHeight() {
        return false;
    }
}
\end{lstlisting}

The remaining \api{Scrollable} methods should still return sensible increments. The example is intentionally incomplete as production code; the point is the width-tracking policy.

A production width-tracking form panel should also answer scroll increments in
terms of the content. A unit increment might be one text row or a fraction of
the default field height. A block increment might be the visible height minus
one row so the user keeps context. These details seem small until keyboard and
wheel scrolling feel either sluggish or disorienting. Layout and navigation are
both part of the same reading experience.

Scrollable forms should keep command rows visible when the workflow demands it.
For a simple dialog, the command row often belongs outside the scroll pane: the
body scrolls, while OK, Cancel, Apply, and Help remain available. For a long
wizard page, the navigation row likewise remains fixed. For a document-like
editor, commands may live in a toolbar outside the document. If the command row
scrolls away, the user may have to hunt for the action that completes the
workflow.

The opposite mistake is to keep too much outside the scroll pane. A fixed
header, fixed warning area, fixed button row, fixed summary, and fixed status
strip can leave a tiny viewport for the actual form. Decide which information
must remain visible during editing. Validation summaries often deserve a fixed
place only when they guide correction; otherwise field-level decoration and a
short status line may be enough.

Scrollable forms also need a policy for initial focus and initial scroll
position. If the first invalid field is below the fold, opening the dialog at
the top may hide the problem. If the dialog always scrolls to the first invalid
field, the user may lose context. Many applications start at the top, then
scroll to the first blocking problem only after OK fails. This policy belongs to
dialog workflow, but the scrollable layout must make it possible.

\textbf{Tables, trees, and split workspaces.}
Not every screen is a form. Many Swing applications are workspaces: a filter bar above a table, a detail pane beside it, a status strip below, perhaps a log or preview panel in another region. MigLayout can express these screens without retreating to nested \api{BorderLayout} panels everywhere.

A common pattern uses one growing row for the main content and one or more fixed rows for filters and status. Columns may split the main area into a growing table and a fixed detail pane. If the detail pane itself contains a form, it can have its own MigLayout with different constraints.

The rule from the previous chapter still applies: use nested panels for genuine geometry changes. A filter bar, main workspace, and status strip have different geometry. They deserve separate panels. Do not use nested panels merely because one parent constraint became hard to understand.

A workspace layout should name the user's primary work surface. In a search
screen, the table may be primary. In an editor, the form may be primary. In a
review screen, the split between list and detail may be primary. Give the
primary surface the main growth claim. Secondary controls should remain
available but should not consume the space that makes the workflow useful. This
principle is often more important than the exact layout manager syntax.

Split panes are sometimes the right tool. MigLayout can place a
\api{JSplitPane}; it should not reimplement user-resizable split behavior with
manual column widths. The same is true for tabbed panes, toolbars, and scroll
panes. Use specialized Swing containers when they own a well-established
interaction. Use MigLayout to arrange those containers into the larger screen.
This keeps layout code from becoming an imitation of components Swing already
provides.

Dense workspaces need restraint. A filter bar with twelve controls may look
powerful but become unreadable. MigLayout can arrange dense controls cleanly,
but density is not the same as usability. Group related filters, provide
advanced sections, use sensible defaults, and let the result table or content
area remain visually dominant. A layout framework can make density survivable;
it cannot decide which controls deserve to be visible all the time.

Status strips should not become junk drawers. A left-aligned status message, a
right-aligned count, and a progress indicator may be useful. Ten labels with
changing widths will cause jitter and distract from work. Reserve stable space
for volatile values, and consider whether some status belongs in the model or
table header instead. MigLayout can align the strip; the screen must decide what
belongs there.

\textbf{Hiding and showing sections.}
Business screens often hide optional sections. MigLayout can react to component visibility when configured appropriately. A hidden advanced-search row should not leave a mysterious blank band. A hidden warning label should not collapse the entire form unpredictably.

When visibility changes affect layout, call \api{revalidate} and \api{repaint} on the relevant panel. The visibility change alters geometry and pixels. If the hidden section lives inside a scroll pane, the viewport and scroll bars may need the new preferred size.

Hidden sections should have an explicit layout mode. Some screens collapse
hidden components completely. Some preserve space to avoid visual jumping. Some
replace a detailed section with a compact summary row. The right answer depends
on workflow. A rarely used advanced-search section should probably collapse. A
validation message area might reserve one line to avoid controls shifting while
the user types. A busy detail area might preserve space so the screen does not
reflow during a load.

When hiding a group, hide the group as a meaningful unit. Do not hide one label
and forget its field. Do not hide the field and leave a required marker behind.
Do not hide the body of an advanced group while its separator remains. Grouping
optional sections into small panels can be legitimate even in MigLayout because
the panel represents a lifecycle and visibility unit, not a workaround for
constraints.

Generated Corusco descriptors can help identify groups, but they should not
decide the visual hide policy. A field may be optional in the model and still
visible in one screen. A field may be read-only in one workflow and hidden in
another. Layout visibility is a presentation decision over model facts.

\textbf{Localization and text growth.}
Layouts that only work in English are not finished. Labels become longer in other languages. Buttons may need wider text. Mnemonics may change. Some components need right-to-left orientation. MigLayout helps because it can let the label column take the width it needs while the field column absorbs the remaining space.

Avoid fixed pixel widths for text-bearing controls unless there is a domain reason. A field for a two-letter country code may be fixed. A customer name field should not be fixed merely because it looked neat in a screenshot. Use column growth and sensible minimum sizes.

Localization also changes grouping. A pair of short English labels may fit
comfortably in one row, while translated labels need two rows. A button strip
may need more width or a different order under platform conventions. A label
with an embedded mnemonic may conflict with another mnemonic after translation.
The layout should tolerate longer text; the resource and command layers should
make the text discoverable; tests or screenshot reviews should include at least
one long-string scenario.

Right-to-left orientation is not only a final polish item. Leading and trailing
alignment, icon placement, table column order, scroll bar placement, and command
row order can all be affected. MigLayout can cooperate with component
orientation when the layout uses relational terms rather than absolute pixel
positions. Code that says "align right" for a trailing command row may need a
platform or orientation review; code that expresses "trailing" as a design idea
is easier to adapt.

Text growth is not restricted to localization. User preferences may increase
font size. Accessibility settings may alter default fonts. FlatLaf themes may
use different metrics. A screen shown through remote desktop may use different
scaling. If a form is built from preferred sizes, growth columns, baseline
alignment, and scrollable bodies, these changes are ordinary inputs. If it is
built from fixed rectangles, each change becomes a bug.

\section{Look-and-Feel and font changes}

Look-and-Feel changes can alter borders, insets, preferred sizes, and baselines. Font changes do the same. A layout that depends on exact preferred sizes from one Look-and-Feel will be brittle. This is another reason to describe relationships instead of coordinates.

The companion examples use \flatlaf{} because it gives consistent modern screenshots, not because the layout should depend on FlatLaf-specific sizes. A well-written MigLayout form should survive reasonable LAF changes with little or no adjustment.

\textbf{Corusco and layout.}
Corusco deliberately does not choose the layout. Generated form companions define fields, models, view contracts, and bindings. The application still builds the panel. This is the right division: generation should not become a GUI builder that hides geometry from the developer.

\begin{coruscobox}{generated forms still need designed panels}
Use generated view interfaces and binding facades to remove repetitive wiring. Use MigLayout to make the screen legible. The two tools solve different problems.
\end{coruscobox}

This division is especially important for teams. Generated code can make field identity stable. It can bind a text field to a form model. It can attach validation decoration. It cannot know whether the field belongs in a compact header, a two-column dialog, a scrollable wizard page, or a dense operations workspace. Human layout design remains necessary.

The boundary is easiest to remember as a question: is this fact about what the
screen means, or where the screen puts it? Field identity, resource keys,
problem targets, help topics, command identity, and dirty state are meaning.
Rows, columns, gaps, spans, split panes, scroll ownership, and command placement
are geometry. Generated code may expose meaning to the view. Handwritten layout
turns that meaning into a readable screen.

This boundary also keeps examples honest. A generated form that constructs its
own panel can be useful as a smoke test, but it is rarely the screen a user
should see. The book's examples should show handwritten MigLayout panels using
generated bindings where appropriate. That demonstrates the actual production
division: Corusco removes repetitive wiring, while the developer still designs
the user's workspace.

There is one exception worth naming. A generator may create a simple diagnostic
or test view so generated models can be exercised without writing a full screen.
Such a view is an aid to testing, not an application layout recommendation. It
should not become the default visual language of a product.

\textbf{A Corusco-shaped form panel.}
The normal Corusco form shape is not a generated panel. It is an application panel that implements a generated view contract. The panel exposes components through methods or fields required by generated bindings. MigLayout arranges those components. A \api{BindingScope} or \api{BehaviorScope} owns the installed relationships.

\begin{lstlisting}[caption={The intended division: handwritten layout, generated binding.}]
CustomerEditFormModel model = new CustomerEditFormModel(customer);
CustomerEditPanel view = new CustomerEditPanel(); // MigLayout inside.
BindingScope scope = new BindingScope();

CustomerEditBindings.install(view, model, scope);
\end{lstlisting}

The generated call should be boring. The layout code should be visible. If the generated binding fails, inspect generated source. If the form looks wrong, inspect MigLayout constraints and component size contracts.

\textbf{Debugging layout.}
Layout debugging begins with preferred sizes, minimum sizes, and grow rules. If a field refuses to expand, inspect the column and component fill behavior. If a panel grows but the scroll pane does not, inspect containment. If a button row jumps, inspect push, wrap, span, and unrelated gaps.

Use temporary borders or background colors when debugging. Give panels names. Print preferred sizes for suspicious components. Reduce the layout to a small reproduction. A layout bug hidden inside a full application is often just a three-component disagreement.

When a component is custom, also inspect the custom component before blaming
the layout. Does \api{getPreferredSize} describe content and insets? Does a
geometry change call \api{revalidate}? Does the component implement
\api{Scrollable} if it is the viewport view? Does it report a useful baseline
if it is meant to align with labels and fields? These questions come directly
from Chapter~\ref{ch:components}; MigLayout cannot infer answers the child does
not provide.

The fastest layout debugging technique is subtraction. Copy the panel into a
small example, remove half the components, and keep the half that still shows
the problem. Repeat until the layout surprise is reduced to a few controls and
one or two constraints. This technique is more useful than staring at a full
screen because many layout bugs are interactions among exactly two facts: a
component preferred size and a missing grow rule, a scroll pane and a
non-tracking view, a hidden component and a missing revalidation, a nested inset
and an outer gap.

MigLayout's own debug mode can be useful during development, but it should not
be a substitute for understanding. Lines and cell markers show where the layout
manager placed things. They do not tell whether the chosen geometry matches the
workflow. Use debug visuals to see the grid, then return to the design
questions: which region is primary, which row grows, which gap expresses a
group, which control should remain visible, and which component is reporting a
surprising size.

Print sizes when visual inspection is ambiguous. Log preferred, minimum,
maximum, and actual sizes for the suspicious component and its parent. Log the
visible rectangle when a scroll pane is involved. Log whether the component is
showing and whether it is visible. These numbers are often enough to reveal that
the component is truthful but the parent is wrong, or the parent is sensible but
the component reports nonsense.

Do not debug layout by adding permanent arbitrary gaps. A top gap of 6 pixels
may hide the immediate overlap, but it also encodes no reason. The next font,
theme, locale, or validation state may break again. If a gap is the right fix,
name the relationship it represents: related, unrelated, paragraph, command
separation, section separation, viewport inset. If no relationship can be
named, the gap is probably compensating for a different defect.

\begin{detailbox}{The question sequence}
When a layout surprises you, ask in order: what is the component's preferred size; which cell contains it; does the column or row grow; does the component fill; does it span; is a scroll pane changing the visible extent; did visibility change require revalidation?
\end{detailbox}

\textbf{When built-in layout managers are enough.}
MigLayout should be the default for serious forms and workspaces, not a reflex for every panel. \api{BorderLayout} is excellent for "header, center, footer" shells. \api{FlowLayout} can be fine for tiny button strips, though MigLayout often gives better alignment. \api{BoxLayout} is useful for simple stacks. Specialized components such as \api{JSplitPane}, \api{JTabbedPane}, and \api{JToolBar} have their own layout behavior.

The pragmatic rule is to use the simplest layout that states the geometry clearly. MigLayout often wins because it states dense geometry clearly. It loses when a standard container already embodies the pattern.

\section{Worked example and review}

\begin{examplebox}{MigLayout}
\migLayoutExample{} builds the form shown in Figure~\ref{fig:miglayout-form}. It is small enough to quote, but it uses the same constraints as larger dialogs: label column, growing field column, preferred rows, bottom push, and right-aligned command row.
\end{examplebox}

\textbf{Case studies and design drills.}
The following small cases are meant to be read as layout diagnoses, not as
separate recipes. Each begins with an unpleasant screenshot symptom and then
asks which part of the geometry contract is missing. This is the practical way
to use MigLayout well: translate visual unease into vocabulary, then adjust the
owner of the relevant relationship.

\textbf{The expanding label.}
A translated label becomes much longer than its English original. The label column expands and steals space from the field column. This may be correct. If the field becomes unusably small, consider wrapping labels, placing labels above fields in narrow contexts, or setting sensible minimum widths for the field column.

\textbf{The notes area that will not grow.}
A \api{JTextArea} sits in a scroll pane, but the notes region remains tiny when the dialog grows. The likely cause is a row that does not grow. The scroll pane may be willing to fill, but no vertical space reaches its row.

\textbf{The mystery horizontal scroll bar.}
A form inside a scroll pane reports a preferred width larger than the viewport and does not track viewport width. The scroll pane dutifully shows a horizontal scroll bar. Fix the view policy; do not merely hide the horizontal scroll bar while leaving clipped content.

\textbf{The vanishing advanced section.}
An advanced section is hidden, but the blank space remains. The layout may be configured to keep hidden components, or the parent was not revalidated. Decide whether hidden sections should collapse. Then make that policy explicit.

\textbf{The generated form that looks generated.}
The team lets generation create too much of the visual structure. The result is a uniform but awkward screen. Move layout back into handwritten panels. Keep generation for descriptors, models, and binding facades.

\textbf{Reviewing the example as a screen.}
The companion form should be reviewed in the same way as a production panel,
only faster. Read the layout constraints aloud before running it: the panel
fills horizontally, has insets, owns a label column and a growing editor column,
keeps ordinary rows at preferred height, pushes the command row toward the
bottom, and aligns commands at the trailing edge. Then resize the window and
ask whether the prediction was true. If the fields grow and the command row
stays stable, the constraint text has earned trust. If the screenshot surprises
the reader, the constraint text needs to be made clearer or the component size
contract needs inspection.

Next perturb the example deliberately. Replace a short label with a long one.
Put the form body into a scroll pane. Add a notes area that should grow
vertically. Hide an optional section and bring it back. Change the Look-and-Feel
font size. These perturbations are not artificial. They are exactly what real
applications do through localization, accessibility, user resizing, validation,
advanced options, and theme changes. A layout that survives them has a design;
a layout that breaks under the first perturbation was a screenshot arrangement.

Finally decide what should become a reusable shape. The two-column form pattern
may deserve a helper in a codebase that uses it everywhere. The exact fields,
groups, and command placement still belong to the screen. A dense filter bar may
deserve a helper for related gaps and trailing commands. A generated form view
contract may deserve a test fixture. The decision should be made after several
screens prove the pattern, not before the first screen is understood. Premature
layout helpers are merely global constraint strings with better names.

This is also where screenshots help without replacing judgement. A screenshot
can show that a dense screen looks balanced under one theme and size. The
constraint review explains why it should remain balanced when the screen is
wider, narrower, translated, or placed in a scroll pane. Keep both forms of
evidence. The picture catches visual awkwardness; the constraint proof catches
fragility before it becomes a visual bug.
Both matter.

\textbf{Checklist for layout.}
\begin{itemize}
\item Prefer MigLayout for nontrivial forms, dialogs, filter bars, and dense workspaces.
\item Keep layout constraints close to the panel they describe.
\item Give exactly the columns and rows that should grow a grow constraint.
\item Distinguish grow from fill, and parent growth from child preferred size.
\item Use nested panels for genuinely different geometry, not to hide confusion.
\item Make scrollable forms track viewport width when horizontal scrolling is not desired.
\item Revalidate and repaint after visibility or preferred-size changes.
\item Treat Look-and-Feel, font, and localization changes as normal layout inputs.
\item Use Corusco for screen facts and bindings; keep geometry readable in application code.
\end{itemize}

MigLayout is also a communication tool for code review. A reviewer should be
able to read the column and row constraints and predict the screen's behavior
under resizing. If the constraints require a screenshot to understand, they are
probably too clever. Prefer a few named layout constants, clear growth columns,
and local comments for unusual constraints over dense strings that only the
author can decode.

The Corusco boundary matters here. Generated field descriptors can say which
field exists, what label resource it uses, which help topic belongs to it, and
which binding should connect it. They should not force a universal visual grid.
Two screens may edit the same record with different emphasis. One may be a
compact dialog; another may be a wide workspace panel. MigLayout lets the view
express that difference while Corusco keeps the facts consistent.

A final review of a MigLayout screen should read like a short design proof. The
screen has one primary region. The primary region receives extra space. The
fixed regions remain useful at preferred size. Scroll ownership is explicit.
Labels and editors align by baseline. Related and unrelated gaps express
meaning. Localization and font growth have a path. Commands are placed by
workflow, not by accident. Corusco bindings attach to named components without
owning geometry. If the proof cannot be stated, the layout is not ready merely
because it looks acceptable today.

This review proof is also what makes the chapter fun rather than fussy.
MigLayout is satisfying when it turns a messy screen into a few clear
sentences. "Two columns; editor grows; notes row consumes height; body scrolls
vertically; buttons stay trailing" is a pleasant thing to read in code. It is
also a pleasant thing for the user, because the screen behaves as if someone
understood the work.

That understanding should remain visible when the screen evolves. A new field
should join an existing row grammar or create a clearly named exception. A new
warning should fit the validation area instead of pushing controls around
unpredictably. A new table should claim height only if it becomes part of the
primary work surface. A new generated binding should attach to an existing
component role, not force the layout to reveal generated structure. Layout
maintenance is therefore a continuation of layout design, not a separate
cleanup phase.

The practical ambition is simple: a future maintainer should be able to change
the screen without fear. They should know which constraint to touch when the
notes field must grow, where to put a new optional field, how to keep the dialog
from scrolling sideways, and why generated companions do not own geometry. That
confidence is the real payoff of using MigLayout as the book's primary layout
framework.

Once this discipline is learned, layout work becomes less theatrical. The team
does not chase pixels; it edits relationships. That is the habit this chapter
tries to make ordinary, and it is also the habit that keeps screenshots stable.
Stable screenshots make review calmer because they are consequences of a stated
geometry contract rather than lucky captures of one window size.

\section{Exercises}

\begin{exercisebox}
\begin{enumerate}
\item Recreate Figure~\ref{fig:miglayout-form} with \api{GridBagLayout}. Compare the amount of code and the review burden.
\item Add a growing notes area above the command row. Predict which row constraint must change before running the example.
\item Place the form in a scroll pane and make it track viewport width but not height.
\item Replace the button text with longer localized strings. Adjust only layout constraints, not component bounds.
\item Implement the same generated-form binding shape with a handwritten MigLayout panel and a Corusco \api{BindingScope}.
\end{enumerate}
\end{exercisebox}
